\documentclass[../../../include/open-logic-section]{subfiles}

\begin{document}

\olfileid{sth}{z}{infinity-again}	
\olsection{无穷}

我们已有的公理足以保证有无穷多个集合存在（前提是至少存在一个集合）。假设存在某个集合，那么根据\olref[sth][z][sep]{prop:emptyexists}，$\emptyset$ 存在。现在，对任意集合~$x$，由\olref[sth][z][pairs]{prop:pairsconsequences}可知集合 $x \cup \{x\}$ 存在。因此，将此操作进行若干次后，我们会得到如下集合：

\begin{enumerate}
	\item[0.] $\emptyset$ %& $0$ & stage $0$\\
	\item[1.] $ \{\emptyset\}$ %& $1$ & stage $1$\\
	\item[2.] $\{\emptyset, \{\emptyset\}\}$ %& $2$ & stage $2$\\
	\item[3.] $\{\emptyset, \{\emptyset\}, \{\emptyset, \{\emptyset\}\}\}$ %&$3$ & stage $3$\\
	\item[4.] $\{\emptyset, \{\emptyset\}, \{\emptyset, \{\emptyset\}\}, \{\emptyset, \{\emptyset\}, \{\emptyset, \{\emptyset\}\}\}\}$% & $4$ & stage $4$ 
\end{enumerate}

并且我们可以验证这些集合两两不同。

我们之所以从 $0$ 开始编号，有几个原因。其中之一是：不难验证我们标记为“$n$”的集合恰好有 $n$ 个成员，并且（直观上）是在第 $n$ 个阶段形成的。

然而，这使得我们拥有\emph{无穷多个}集合，但并不能保证存在一个\emph{无穷集合}，即一个拥有无穷多成员的集合。而这确实很重要：除非我们能找到一个（戴德金）无穷集合，否则我们就无法构造一个戴德金代数。但我们需要一个戴德金代数，以便将其视为自然数集。（比较\olref[sfr][infinite][dedekindsproof]{sec}。）

重要的是，我们迄今提出的公理\emph{并不}保证存在任何无穷集合。因此，我们必须提出一条新公理：

\begin{axiom}[无穷公理]
存在一个集合 $I$，使得 $\emptyset \in I$，并且只要 $x \in I$，就有 $x \cup \{x\} \in I$。
\begin{align*}
	\exists I( & (\exists o \in I)\forall x\ x \notin o \land {}\\
	& (\forall x \in I)(\exists s \in I)\forall z(z \in s \liff (z \in x \lor z = x)))
\end{align*}
\end{axiom}

容易看出，无穷公理提供的集合 $I$ 是戴德金无穷的。它的特异元素是 $\emptyset$，其上的单射由 $s(x) = x\cup
\{x\}$ 给出。现在，\olref[sfr][infinite][dedekind]{thm:DedekindInfiniteAlgebra} 展示了如何从一个戴德金无穷集合中提取出一个戴德金代数；我们将把它视为我们的自然数集。更精确地说：

\begin{defn}\ollabel{defnomega}
设 $I$ 为无穷公理提供的任意集合。令 $s$ 为函数 $s(x) = x \cup \{x\}$。令 $\omega = \closureofunder{s}{\emptyset}$。我们称 $\omega$ 的成员为\emph{自然数}，并称 $n$ 是将 $s$ 作用于 $\emptyset$ $n$ 次的结果。
\end{defn}

现在你可以回顾并验证，前面几段中标记为“$n$”的集合将被\emph{作为}数 $n$ 来对待。

我们将在\olref[sth][z][nat]{sec}讨论这一规定的意义。现在，它使我们能够证明一个直观的结果：

\begin{prop}\ollabel{naturalnumbersarentinfinite}
任何自然数都不是戴德金无穷的。
\end{prop}

\begin{proof}
证明采用归纳法，即\olref[sfr][infinite][induction]{thm:dedinfiniteinduction}。显然 $0=\emptyset$ 不是戴德金无穷的。对于归纳步骤，我们将证明其逆否命题：如果（荒谬地）$s(n)$ 是戴德金无穷的，那么 $n$ 也是戴德金无穷的。

因此，假设 $s(n)$ 是戴德金无穷的，即存在某个单射 $f$ 满足 $\ran{f}\subsetneq \dom{f} = s(n) = n \cup
\{n\}$。需要考虑两种情况。

\emph{情况 1：$n \notin \ran{f}$。} 所以 $\ran{f} \subseteq n$，且 $f(n) \in n$。令 $g = \funrestrictionto{f}{n}$；则 $\ran{g} =
\ran{f} \setminus \{f(n)\} \subsetneq n = \dom{g}$。因此 $n$ 是戴德金无穷的。

\emph{情况 2：$n \in \ran{f}$。} 取定 $m \in \dom{f} \setminus \ran{f}$，并定义一个定义域为 $s(n) = n \cup \{n\}$ 的函数 $h$：
\[
h(x) = 
	\begin{cases}
		f(x) & \text{若 }f(x) \neq n\\
		m & \text{若 }f(x)=n
	\end{cases}
\]
所以，除了 $h(f^{-1}(n)) = m
\neq n = f(f^{-1}(n))$ 之外，$h$ 和 $f$ 处处一致。由于 $f$ 是单射，$n \notin \ran{h}$；并且 $\ran{h} \subsetneq \dom{h} = s(n)$。现在利用情况 1 的论证，可知 $n$ 是戴德金无穷的。
\end{proof}

然而，问题依然存在：我们该如何\emph{辩护}无穷公理？简短的回答是，我们需要向我们一直讲述的故事中引入另一条原则。该原则如下：

\begin{enumerate}
	\item[] \stagesinf. 存在一个无穷阶段。也就是说，存在一个阶段，它 (a) 不是第一阶段，(b) 前面有其他阶段，但 (c) 没有直接前驱。
\end{enumerate}

无穷公理可直接由这一原则得出。
我们知道，自然数 $n$ 在第 $n$ 个阶段形成。
所以集合 $\omega$ 在第一个无穷阶段形成。
而 $\omega$ 本身就见证了无穷公理。

然而，这仅仅是将问题又推了回去：我们如何能为 \stagesinf 辩护呢？
与 \stagessucc 一样，它并非我们关于累积迭代层次所讲述故事的显式部分。
不仅如此：在集合逐阶段形成的迭代层次这一观念本身中，没有任何东西迫使我们认为该过程包含一个\emph{无穷}阶段。
认为这些阶段的排序如同自然数一样，似乎是完全自洽的。

然而，这引出了一个明显的问题。
在 \olref[sfr][infinite][dedekindsproof]{sec} 中，我们考察了戴德金关于存在一个（由思想构成的）戴德金无穷集合的“证明”。
这可能并未让你觉得十分令人满意。
但如果 \stagesinf 并非由集合的迭代观念（或“思维法则”）“强加给我们的”，那么对于存在戴德金无穷集合的主张，我们仍然缺乏内在的辩护。

当然，关于这点还有更多可说的。
但希望你现已能够开始思考，辩护一条公理（或原则）究竟需要\emph{什么}。
在接下来的内容中，我们将直接假定 \stagesinf{} 成立。

\end{document}